\documentclass[10pt,aspectratio=169]{beamer}
\usetheme{metropolis}
\usepackage[T1]{fontenc}
\usepackage[utf8]{inputenc}
\usepackage{textcomp}
\usepackage{booktabs}
\usepackage{multirow}
\usepackage{amsmath,amssymb}
\usepackage{graphicx}
\usepackage{array}
\setbeamertemplate{footline}[frame number]
\title{Electoral Turnover and Forest Cover in Colombia}
\subtitle{Mayoral elections 2011--2023, MOD44B VCF tree cover 2000--2025}
\author{Anzony Quispe Rojas}
\date{\today}
\begin{document}\begin{frame}\titlepage\end{frame}

\begin{frame}{Question and approach}
\begin{itemize}
  \item \textbf{Question.} Do mayoral electoral turnovers in Colombia change
        forest cover in the years following the election?
  \item \textbf{Setting.} 1{,}099 municipalities with mayoral elections in
        2011, 2015, 2019, 2023 (4-year terms).
  \item \textbf{Strategy.} Adapt Marx, Pons \& Rollet (2024) outcome
        construction; estimate via 2SLS with rank-based instruments rather
        than a close-elections RDD.
  \item \textbf{Treatment.} $T_{m,t}=1$ if the winning party in election
        $t$ differs from the winning party in $t-1$.
  \item \textbf{Instruments.}
      $Z = \mathbf{1}\{\text{previous winner-party rank in election } t > 1\}$ (user-specified, sharp);
      $Z_2 = \mathbf{1}\{\text{previous runner-up wins election } t\}$ (fuzzy alternative).
\end{itemize}
\end{frame}


\begin{frame}{Data sources}
\begin{itemize}
  \item \textbf{Elections.} \texttt{elections\_long\_clean\_with\_details.dta}
        (19{,}597 candidate rows, 1{,}102 munis, 2011/15/19/23). Collapsed
        to 3,287 muni $\times$ election cells where $T$ is defined.
  \item \textbf{Boundaries.} DANE \emph{Marco Geoest{\'a}distico Nacional}
        municipal layer (\texttt{MGN\_MPIO\_POLITICO}, EPSG:4686), 1{,}121 munis.
  \item \textbf{Forest cover.} User-supplied GEE export of MOD44B v061 VCF
        \emph{Percent\_Tree\_Cover} averaged inside each muni polygon:
        \texttt{VCF\_TreeCover\_Municipios\_Colombia\_2000\_2025.csv}
        (2000--2025, 26 years; 1{,}121 munis, 29{,}146 cells, no missingness).
\end{itemize}
\end{frame}


\begin{frame}{Treatment and instruments}
\begin{block}{Treatment}
$T_{m,t} = \mathbf{1}\{\,\text{winner party}_{m,t} \neq
\text{winner party}_{m,t-1}\,\}$.
Across the analysis sample, $\bar T = 0.854$ (turnover is the rule).
\end{block}
\begin{block}{Instrument $Z$ (user spec, ``rank as instrument'')}
$Z_{m,t}=\mathbf{1}\{\,\text{rank of the previous-period winning party
in the current election}>1\,\}$, defined when that party re-ran.
\end{block}
\begin{block}{Note}
On the eligible subsample (incumbent party re-ran), $Z\equiv T$ by
construction. The 2SLS coefficient on $T$ therefore equals OLS on that
subsample. We also report a truly fuzzy alternative $Z_2=\mathbf{1}\{
\text{previous runner-up wins}\}$.
\end{block}
\end{frame}


\begin{frame}{Outcome construction (Marx--Pons--Rollet \S3.3)}
For an election in muni $m$ at year $t$ and outcome $Y$ (mean
\emph{Percent\_Tree\_Cover} inside the muni polygon):
\begin{align*}
\Delta Y^{\text{main}}_{m,t} &= \tfrac{1}{4}\!\sum_{\tau=1}^{4} Y_{m,t+\tau} - Y_{m,t-1}\\
\Delta Y^{k}_{m,t} &= \tfrac{1}{k}\!\sum_{\tau=1}^{k} Y_{m,t+\tau} - Y_{m,t-1}, \quad k\in\{1,2,3\}\\
\Delta Y^{\text{pre}_3}_{m,t} &= \tfrac{1}{4}\!\sum_{\tau=1}^{4} Y_{m,t+\tau} - \tfrac{1}{3}\!\sum_{\tau=1}^{3} Y_{m,t-\tau}
\end{align*}
Each $\Delta Y$ is z-standardised so coefficients are in SD units. Higher tree
cover is ``good''; no sign flip.
\end{frame}


\begin{frame}{Sample splits by baseline forest cover (year 2000)}
\begin{columns}[T,onlytextwidth]
\column{.52\textwidth}
\centering
\includegraphics[width=\linewidth]{/workspace/A_MicroData/figs/fig_baseline_hist.png}\\[-2pt]
\footnotesize Distribution of \emph{Percent\_Tree\_Cover} across 1{,}099
munis in 2000.
\column{.46\textwidth}
\begin{itemize}
  \item Median baseline = \textbf{30.56\%}.
  \item \textbf{Above}: 1,642 muni-elections.
  \item \textbf{Below}: 1,645 muni-elections.
  \item Same regressions are run on \texttt{all}, \texttt{above}, and
        \texttt{below} samples.
\end{itemize}
\end{columns}
\end{frame}


\begin{frame}{Geography of the baseline split}
\centering
\includegraphics[height=.78\textheight]{/workspace/A_MicroData/figs/fig_baseline_split.png}
\end{frame}


\begin{frame}{Forest-cover trend, national mean}
\centering
\includegraphics[height=.7\textheight]{/workspace/A_MicroData/figs/fig_treecover_trend.png}
\end{frame}


\begin{frame}{Descriptives: turnover and victory margin}
\centering
\includegraphics[height=.4\textheight]{/workspace/A_MicroData/figs/fig_turnover_rate.png}
\quad
\includegraphics[height=.4\textheight]{/workspace/A_MicroData/figs/fig_margin_hist.png}\\
\small Left: turnover rate by election year. Right: distribution of
$\text{share}_\text{winner} - \text{share}_\text{runner-up}$.
\end{frame}


\begin{frame}{Geography of turnover, 2023}
\centering
\includegraphics[height=.78\textheight]{/workspace/A_MicroData/figs/fig_map_turnover_2023.png}
\end{frame}


\begin{frame}{First stages}
\centering
\renewcommand{\arraystretch}{1.15}\footnotesize
\begin{tabular}{llcccc}
\toprule
Sample & Instrument & Coef. & SE & $p$ & $N$ \\
\midrule
all & Z (incumbent rank \textgreater{} 1 \textbar{} re-ran) & 1.000$^{***}$ & 0.000 & 0.000 & 1,517 \\
all & Z2 (prev runner-up wins) & 0.218$^{***}$ & 0.010 & 0.000 & 3,276 \\
above & Z (incumbent rank \textgreater{} 1 \textbar{} re-ran) & 1.000$^{***}$ & 0.000 & 0.000 & 769 \\
above & Z2 (prev runner-up wins) & 0.217$^{***}$ & 0.015 & 0.000 & 1,634 \\
below & Z (incumbent rank \textgreater{} 1 \textbar{} re-ran) & 1.000$^{***}$ & 0.000 & 0.000 & 748 \\
below & Z2 (prev runner-up wins) & 0.223$^{***}$ & 0.015 & 0.000 & 1,642 \\
\bottomrule
\end{tabular}
\\[0.5em]
\footnotesize SE clustered at \texttt{mpio}; two-way FE on
\texttt{coddepto} and \texttt{election\_year}. $Z$ has SE $=0$ because
of the mechanical relationship $Z\equiv T$ on the eligible subsample.
\end{frame}


\begin{frame}{Main results — whole sample}
\centering
\renewcommand{\arraystretch}{1.15}\footnotesize
\begin{tabular}{lccccc}
\toprule
Specification & Main & $k{=}1$ & $k{=}2$ & $k{=}3$ & Pre$_3$ \\
\midrule
OLS (FE) & -0.018 & -0.019 & -0.035 & -0.030 & -0.025 \\
 & (0.041) & (0.046) & (0.039) & (0.041) & (0.042) \\
$N$ & 3,287 & 3,287 & 3,287 & 3,287 & 3,287 \\
\addlinespace[2pt]
2SLS, $Z=\mathbf{1}\{\text{rank}_\text{inc}{>}1\}$ & -0.005 & -0.004 & -0.031 & -0.015 & -0.023 \\
 & (0.046) & (0.050) & (0.043) & (0.046) & (0.048) \\
$N$ & 1,517 & 1,517 & 1,517 & 1,517 & 1,517 \\
\addlinespace[2pt]
2SLS, $Z_2=\mathbf{1}\{\text{prev RU wins}\}$ & -0.094 & 0.113 & -0.060 & -0.084 & -0.139 \\
 & (0.185) & (0.187) & (0.175) & (0.184) & (0.183) \\
$N$ & 3,276 & 3,276 & 3,276 & 3,276 & 3,276 \\
\addlinespace[2pt]
\bottomrule
\end{tabular}
\\[0.3em]
\footnotesize Outcomes are z-scored. SE in parentheses, clustered at
\texttt{mpio}. FE on \texttt{coddepto} $\times$ \texttt{election\_year}.
\end{frame}


\begin{frame}{Results — above-median 2000 tree cover}
\centering
\renewcommand{\arraystretch}{1.15}\footnotesize
\begin{tabular}{lccccc}
\toprule
Specification & Main & $k{=}1$ & $k{=}2$ & $k{=}3$ & Pre$_3$ \\
\midrule
OLS (FE) & -0.009 & -0.031 & -0.050 & -0.030 & 0.017 \\
 & (0.063) & (0.065) & (0.058) & (0.063) & (0.065) \\
$N$ & 1,642 & 1,642 & 1,642 & 1,642 & 1,642 \\
\addlinespace[2pt]
2SLS, $Z=\mathbf{1}\{\text{rank}_\text{inc}{>}1\}$ & -0.041 & -0.038 & -0.078 & -0.059 & -0.031 \\
 & (0.068) & (0.071) & (0.062) & (0.068) & (0.070) \\
$N$ & 769 & 769 & 769 & 769 & 769 \\
\addlinespace[2pt]
2SLS, $Z_2=\mathbf{1}\{\text{prev RU wins}\}$ & -0.199 & -0.035 & -0.203 & -0.153 & -0.134 \\
 & (0.259) & (0.275) & (0.238) & (0.256) & (0.262) \\
$N$ & 1,634 & 1,634 & 1,634 & 1,634 & 1,634 \\
\addlinespace[2pt]
\bottomrule
\end{tabular}
\\[0.3em]
\footnotesize Munis with baseline tree cover above 30.6\% in 2000.
\end{frame}


\begin{frame}{Results — below-median 2000 tree cover}
\centering
\renewcommand{\arraystretch}{1.15}\footnotesize
\begin{tabular}{lccccc}
\toprule
Specification & Main & $k{=}1$ & $k{=}2$ & $k{=}3$ & Pre$_3$ \\
\midrule
OLS (FE) & -0.024 & 0.001 & -0.021 & -0.030 & -0.066 \\
 & (0.054) & (0.066) & (0.053) & (0.055) & (0.053) \\
$N$ & 1,645 & 1,645 & 1,645 & 1,645 & 1,645 \\
\addlinespace[2pt]
2SLS, $Z=\mathbf{1}\{\text{rank}_\text{inc}{>}1\}$ & 0.042 & 0.051 & 0.027 & 0.040 & -0.018 \\
 & (0.062) & (0.070) & (0.061) & (0.063) & (0.066) \\
$N$ & 748 & 748 & 748 & 748 & 748 \\
\addlinespace[2pt]
2SLS, $Z_2=\mathbf{1}\{\text{prev RU wins}\}$ & 0.035 & 0.297 & 0.092 & 0.002 & -0.125 \\
 & (0.261) & (0.252) & (0.250) & (0.260) & (0.257) \\
$N$ & 1,642 & 1,642 & 1,642 & 1,642 & 1,642 \\
\addlinespace[2pt]
\bottomrule
\end{tabular}
\\[0.3em]
\footnotesize Munis with baseline tree cover at or below 30.6\% in 2000.
\end{frame}


\begin{frame}{Take-aways}
\begin{itemize}
  \item Turnover rate is extraordinarily high in Colombian mayoral elections
        (\textbf{$\bar T \approx 0.85$}). 2023 alone has 93\% turnover.
  \item Across all samples, specifications, and instruments, the effect of
        turnover on $\Delta$ tree cover is small (mostly $|\hat\beta|<0.10$ SD)
        and statistically indistinguishable from zero.
  \item Signs are slightly negative in the whole and above-median samples
        and mixed in the below-median sample; none survive clustered SEs.
  \item The user-specified instrument $Z$ is mechanically equal to $T$
        on the eligible subsample (sharp design), so its 2SLS coefficient
        merely re-states the OLS coefficient. The fuzzy $Z_2$ recovers a
        runner-up-becomes-mayor LATE; estimates are noisier but consistent
        in sign with OLS.
\end{itemize}
\end{frame}


\begin{frame}{Caveats and next steps}
\begin{itemize}
  \item $T$ is undefined for the 2011 election (no $t{-}1$ winner in the
        data); the analysis sample is therefore 2015, 2019, 2023.
  \item For 2023, \texttt{partyid} is missing for $\sim$60\% of candidates,
        so the cross-year party match uses the cleaned party \emph{name}
        string. Party-name churn between elections could induce
        measurement error in $T$.
  \item Forest cover from MOD44B VCF is a slow-moving, coarse (250 m)
        outcome. A more responsive outcome (Hansen GFC annual loss,
        active fires) might pick up policy changes that VCF smooths over.
  \item No covariates beyond fixed effects; ideology, coalition status,
        and incumbent characteristics are available in the panel for
        future robustness checks.
\end{itemize}
\end{frame}


\begin{frame}{Files produced}
\footnotesize
\begin{tabular}{ll}
\toprule
File & Contents \\
\midrule
\texttt{A\_MicroData/analysis\_panel.dta} & Master Stata file: T, Z, Z2, $\Delta Y$, baseline\_2000 \\
\texttt{A\_MicroData/analysis\_panel.parquet} & Same panel in parquet \\
\texttt{A\_MicroData/elections\_panel.parquet} & Pre-merge elections panel \\
\texttt{A\_MicroData/baseline\_tree\_2000.csv} & 1{,}121 munis × year-2000 tree cover \\
\texttt{A\_MicroData/results/iv\_\{all,above,below\}.csv} & Python results, by sample \\
\texttt{A\_MicroData/results/first\_stage.csv} & First-stage coefficients \\
\texttt{C\_Programs/run\_analysis.do} & \textbf{Stata do-file (deliverable)} \\
\texttt{C\_Programs/build\_*.py} & Python pipeline \\
\texttt{C\_Programs/slides.tex / .pdf} & This deck \\
\bottomrule
\end{tabular}
\end{frame}
\end{document}
